% ID: SEL-D0008-BD-E08
% ESTADO: cotejada. Segunda pasada visual completa y borrador compilado revisados.
% Fuente: Selectividad/matematicas_aplicadas/ORD_SUPLENTE_Modelo-4-Enunciados (Mat-Ap-2021).pdf
% SHA-256: 74b067b81be20aa0262a0cc58795f3b25488aa9c749d825b7b33f656dd1ae089
% Página: 2; bloque D; ejercicio original 8.

El peso de los paquetes de arroz de una marca comercial concreta sigue
una ley Normal de media $1000\,g$ y varianza $256\,g^2$.
\begin{enumerate}[label=\alph*)]
  \item \textbf{(0.75 puntos)} Calcule la probabilidad de que el peso medio
  de las muestras de tamaño 64 sea menor que $996\,g$.
  \item \textbf{(1.5 puntos)} Tras varias denuncias presentadas por falta de
  peso en los citados paquetes, una organización de consumidores ha procedido
  a tomar una muestra de 64 paquetes, resultando que la suma de los pesos ha
  sido de $63\,744\,g$. Halle un intervalo de confianza al 90\% para estimar
  el peso medio real de los paquetes de arroz de esa marca.
  \item \textbf{(0.25 puntos)} A la vista del intervalo obtenido y teniendo
  en cuenta que el peso que marca el paquete es de $1000\,g$,
  \textquestiondown{}cree que la denuncia tiene base?
\end{enumerate}
